 %%%%%%%%%%%%%%%%%%%%%%%%%%%%%%%%%%%%%%%%%% 
 % @File    : c:\Users\Administrator\Desktop\Econometrics\太极图\TaiJi.tex
 % @Date    : 2021-02-10 11:03:26
 % @Author  : RankFan
 % @Email   : 1917703489@qq.com
 % -----
 % Last Modified: 2021-02-10 11:05:53
 % Modified By: Rank_fan
 % -----
 %%%%%%%%%%%%%%%%%%%%%%%%%%%%%%%%%%%%%%%%%% 
 \documentclass[border=1cm, 12pt]{standalone}
\usepackage{tikz}
\begin{document}
    \begin{tikzpicture}
        \draw(0,0) circle(3cm);  %画一个4cm的圆圈
        %把圆的右半边填黑
        \begin{scope}
            \clip(0,0) circle(3cm);
            \fill[black] (0,-3) rectangle (3,3);
        \end{scope}
        %填黑八卦图左边的半圆
        \begin{scope}
            \clip(0,1.5) circle(1.5cm);
            \fill[white] (-3,0) rectangle (3,3);
        \end{scope}
        %填白八卦图右边的半圆
        \begin{scope}
            \clip(0,-1.5) circle(1.5cm);
            \fill[black] (-3,0) rectangle (3,-3);
        \end{scope}
        %把黑色部分的小圆圈填为白色
        \begin{scope}
            \clip(0,1.5) circle(0.5cm);
            \fill[black] (-3,0) rectangle (3,3);
        \end{scope}
    %绘制下面的白色圆圈
        \begin{scope}
            \clip(0,-1.5) circle(0.5cm);
            \fill[white] (-3,0) rectangle (3,-3);
        \end{scope}
        % \node[elegant,black](a) at (0,3.5) {——}
    \end{tikzpicture}	
\end{document}